\chapter{Cross-Contribution Discussion}
\label{chap:discussion}

% Drafted under T-004 from the current text of Chapters 3 to 6, and revised on 11 September 2026
% from the author's proposed text (tmp/chapter7-review/chapter7-proposed.md). Constraints from
% notes/dissertation-structure-v0.3.md: the comparison is conceptual, no result of Hybroid is
% compared numerically with a result of HGANN-Mal, and no developmental sequence from the IoT
% work to the Android work is claimed. Section 5.3 cites sec:discussion:3 by label.
% Length fixed at four pages by D-044.

The contribution chapters examine different forms of context under separate experimental
protocols. This discussion relates the observations each representation requires to the
conclusions its evaluation supports. Comparisons between \hybroid{} and \hgann{} are conceptual;
their scores do not provide a direct comparison because the corpora and tasks differ. The central
question is when additional context provides useful information and how the evidence for that
benefit is bounded.

\section{Communication Context and Program Context}
\label{sec:discussion:1}

In the network studies of Chapter~\ref{chap:anomaly}, context comprises the communication
partners of a device or service and the exchanges observed between them. In \hgann{}, context
consists of relations among methods within an application. \hybroid{} combines program structure
with network observations. Each representation depends on its observation process: the service
graph contains relations encountered during learning, whereas the call graph contains
invocations recovered by static analysis. Neither is a complete record of system behaviour. This
comparison concerns their methodological properties; it does not imply a developmental link
between the IoT and Android studies.

Legitimate changes can also produce anomalous observations. Each of the nine false alarm episodes
on the industrial benign monitoring day had a benign cause. Three followed engineering
workstation activity and two followed planned device restarts. The service-graph design
permits an operator to approve new relations. An application update that changes its call graph
presents a related problem, but the Android evaluations did not examine it. The industrial
observations therefore motivate a question about adaptation to legitimate change, without
demonstrating its effect on either Android detector.

\section{Multimodal and Higher-Order Representations}
\label{sec:discussion:2}

Both Android systems derive method-level call graphs with Androguard and use opcode information
with operands discarded. \hybroid{} adds traffic captured from the application. \hgann{}
constructs hyperedges over information derived from its static representation. The additions
address different observational limits. Flow averaging removes the order and detailed temporal
organisation of exchanges, while retaining aggregate features such as mean record duration.
Hypergraph construction remains dependent on the calls and methods recovered by the analyser.

Within the \hybroid{} evaluation, fusion increased detection $F_1$ by one or two percentage
points over the better single modality. For category classification, its effect ranged from a
two-point decrease to a three-point increase. Network features led code features for detection
and trailed them for category classification under each reported learner. Thus, the observed
benefit of an additional modality depends on the task and classifier.

The hypergraph systems in Chapter~\ref{chap:hgann} exceeded the pairwise baselines under the
reported protocols. \hgann{}'s advantage over the non-attentive systems was metric-dependent: on
CICMalDroid binary detection, it trailed HGNN+ by 0.1 accuracy points but led by 0.3 $F_1$
points. Hyperedge construction and attention changed together, so these differences cannot be
attributed to attention alone. The results support comparisons between complete systems; the
missing ablation limits conclusions about the mechanism responsible.

\section{Observation Requirements and Data-Acquisition Cost}
\label{sec:discussion:3}

The fused \hybroid{} model requires a traffic capture for each application. Its evaluation used
CICAndMal2017 captures recorded by the dataset authors on three physical handsets under scripted
use. Both branches excluded 47 graph-extraction failures from the 2,126 matched samples, leaving
2,079 for evaluation. The static branch was also evaluated on 135,905 applications without
corresponding traffic records. This demonstrates its use on a larger corpus, although it does not
quantify a runtime advantage over the fused model.
\hgann{} likewise requires application packages without traffic acquisition, but reports no
extraction-time or memory measurements.

The industrial study required access to the plant and manual labelling of 25 attack episodes.
The Raspberry Pi~3 experiment measured extractor throughput separately, including a moderate
load of about 2,500 input frames per second. Collection effort and extraction throughput are
therefore different costs. The available records do not isolate how either determined corpus
size.

These requirements affect which applications and operating conditions can be evaluated. The
S7-300 test partition contains one episode per attack class; each \hybroid{} malware category
contains between 101 and 112 applications in the corpus. Additional traffic may provide evidence
absent from static code, but its benefit must justify the required observation.
Chapter~\ref{chap:hybroid}'s hypothesis that unreachable injected calls could leave traffic
unchanged remains untested against \hybroid{}.

\section{Predictive Evaluation and Adversarial Evaluation}
\label{sec:discussion:4}

The predictive results of Chapters~\ref{chap:anomaly}--\ref{chap:hgann} concern samples and
attacks that were not adapted to the evaluated detector. Chapter~\ref{chap:advml} adds an
explicit adversary and shows why attack performance must be interpreted alongside useful
classification. A constant classifier has zero accuracy loss under every label-preserving attack
(Equation~\ref{eq:advml:constant}). One SD configuration reportedly reached constant attribute
predictions at approximately 80 per cent accuracy. This makes a matched trivial baseline relevant
to interpretation, even when an aggregate score appears high.

For Android malware, the attacker must also preserve a functioning application and its malicious
behaviour. Published attacks can satisfy such constraints while changing recovered call graphs
\cite{song2026fcghunter}. Both Android systems derive their program representations from those
graphs; constructing hyperedges does not remove the upstream influence of code modification. An
evaluation would need to specify valid package transformations and attack the complete pipeline,
including the fusion stage where applicable. Clean and attacked performance would then be
reported under the same conditions \cite{carlini2019evaluating}. Neither Android system underwent
that evaluation, so their predictive results establish neither adversarial robustness nor a
demonstrated exploitable weakness.

\section{Dataset Limitations and Extraction Failures}
\label{sec:discussion:5}

Evaluation depends on which samples reach the classifier. \hybroid{} excluded the same 47
graph-extraction failures from both branches. \hgann{} does not report its extraction-failure count or how failures were handled. Their
extraction coverage therefore cannot be compared. CICMalDroid retains 11,598 of 17,341 collected applications after its
analysis filters, which may exclude some evasive behaviour. The industrial extractor requires
observed traffic, while the hidden contest procedure excludes images that fail its
face-detection check. Each result is conditional on the samples retained by its pipeline.

Application provenance and traffic acquisition also require separate treatment. CICAndMal2017
collected benign applications from Google Play and malware from VirusTotal and Contagio, although
their traffic was captured using common handsets and procedures. Shared capture conditions
control part of the measurement process; they do not eliminate differences in application
origin. The Drebin experiment combines malware with benign applications drawn from AndroZoo, with
no documented collection-window restriction on that draw. Such differences can provide
predictive cues unrelated to malicious behaviour, as demonstrated in the Android case study of
Arp et al.~\cite{arp2022dos}.

The contribution chapters identify inconsistencies in their source records and recompute the
quantities supported by the available counts. Remaining unknowns limit interpretation even where
a reported score can be reproduced arithmetically.

\section{Scope of Generalisation}
\label{sec:discussion:6}

Table~\ref{tab:discussion:scope} distinguishes evaluation protocols and the populations behind
the reported class shares. These percentages describe binary detection; they do not characterise
the multiclass label distributions. Corpus proportions need not equal an undocumented
test partition.

\begin{table}[tbp]
  \caption[Evaluation conditions and reported binary class shares]{Evaluation conditions and
  reported binary class shares. ``Unreported'' denotes missing information, rather than a
  measured absence. Baseline descriptions refer to the evidence presented in the contribution
  chapters.}
  \label{tab:discussion:scope}
  \centering
  \small
  \begin{tabularx}{\textwidth}{@{}>{\raggedright\arraybackslash}p{2.9cm}
    *{4}{>{\raggedright\arraybackslash}X}@{}}
    \toprule
    Study & Reported split & Positive class and share & Population represented & Trivial
      baseline \\
    \midrule
    S7-300 & Whole capture sessions & Attack, 6.2\% & Confirmed test records & Stated in the
      text \\
    SWaT GAN & Chronological; later 393,679 timestamps for testing & Anomalous, 13.1\% &
      Later part of one SWaT attack recording & Stated in the text \\
    \hybroid{} & Five-fold cross-validation & Malware, 20.0\% in source corpus & 2,079 of 2,126
      samples & Not reported \\
    \hgann{}: Drebin/AndroZoo & Random 70/20/10 & Malware, approximately 37.0\% & Assembled
      binary corpus & Corpus majority \\
    \hgann{}: CICMalDroid & Random 70/20/10 & Malware, 84.5\% & Full study corpus & Corpus
      majority \\
    SAFAIR histopathology benchmark & Evaluated subset unreported & IDC, unreported & Test
      population unidentified & Analytical discussion only \\
    \bottomrule
  \end{tabularx}
\end{table}

The SWaT protocol evaluates the later 393,679 timestamps of the same attack-period recording;
the earlier 56,240 timestamps support training and validation. This split prevents reuse of test
observations during model development, but it does not measure transfer to another campaign or
installation. Session separation and random partitioning address different questions from a
time-ordered evaluation \cite{pendlebury2019tesseract}. The reported folds also do not provide
the independent repeated-run outcomes needed to quantify variability. In the industrial
studies, the S7-300 test attacks belong to classes present during training, and the SWaT detector
is developed with attack observations in training and supervised selection.

Within-study contrasts provide the most direct evidence for performance differences under their
reported conditions. Their interpretation still depends on the controls: \hybroid{} does not
specify whether preprocessing and embeddings were fitted within each fold, and \hgann{} does not
isolate attention from hyperedge construction. These limitations concern internal validity.
Transfer to another population is a separate question, as illustrated by performance
deterioration under challenging conditions in the comparative Android study of Gao
et al.~\cite{gao2024comprehensive}. Chapter~\ref{chap:advml}'s analytical identities apply
independently of a particular corpus under their stated assumptions; its empirical measurements
retain their sampling and implementation limits.

\section{Implications for Android Malware Analysis}
\label{sec:discussion:7}

The results support choosing context according to the prediction task and the observations
available. Traffic improved detection relative to code features in the \hybroid{} experiment, but
fusion did not improve every category-classification configuration. Higher-order systems
performed well in the Chapter~\ref{chap:hgann} comparisons, although the contribution of
attention remains unidentified. Additional context therefore requires evidence of benefit within
the intended task, together with controls that identify what produced the difference.

Observation limits remain after a representation is enriched. Static analysis avoids
per-application traffic acquisition but cannot recover behaviour absent from its analysis.
Traffic depends on the exercised behaviour and capture environment; the protocol mix in the
\hybroid{} corpus further limits conclusions about encrypted communication. A deployment claim
would require evaluation of these conditions as well as prediction quality. The dissertation
establishes bounded benefits of contextual representations and a method for examining the
validity of their evaluation. Its remaining empirical questions concern the missing component
ablations and performance under valid application transformations or later data.
